%% ============================================================ 6 SAMPLING AND THE CENTRAL LIMIT THEOREM
\clearpage
\inchap{Sampling and the Central Limit Theorem}
\section{Population and Sample}
\begin{center}\ba\small
\begin{tabular}{|M{3.8cm}|L{7.2cm}|L{4.2cm}|}\hline
\textbf{Term} & \textbf{Meaning} & \textbf{Notation}\\\Xhline{2pt}
Population & the whole group of interest & mean $\mu$, variance $\sigma^{2}$ (parameters)\\\hline
Sample & a part of the population, size $n$ & $\xb$, $s^{2}$ (statistics)\\\hline
Random sample & \begin{itemize}[leftmargin=*,nosep,before=\vspace{3pt}]\item every member of the population is equally likely to be chosen\item the selection is independent\end{itemize}\textbf{Benefits:} avoid bias, and ensure the sample is representative of the population.\newline \textbf{How to obtain:} number every member of the population with a distinct positive integer from $1$ to $N$. Use a computer to randomly generate $n$ distinct positive integers. Select the members whose numbers correspond to those generated by the computer. & \\\hline
\end{tabular}\end{center}
\begin{center}\ba\small
\begin{tabular}{|L{7.2cm}|L{8.0cm}|}\hline
\textbf{Not random if \ldots} & \textbf{Reason}\\\Xhline{2pt}
chosen by convenience (e.g.\ first 30 in the canteen, volunteers) & some members are more likely than others\\\hline
only part of the population can be chosen & the rest has no chance\\\hline
one choice affects another (e.g.\ friends chosen together) & not independent\\\hline
\end{tabular}\end{center}
\section{Unbiased Estimates}
\begin{center}\bba\small
\begin{tabular}{|M{3.4cm}|M{6.0cm}|M{5.8cm}|}\hline
\textbf{Estimate of} & \textbf{Raw data} & \textbf{Coded: $\Sigma(x-a)$, $\Sigma(x-a)^{2}$}\\\Xhline{2pt}
Mean $\mu$ & $\xb=\dfrac{\sum x}{n}$ & $\xb=a+\dfrac{\sum(x-a)}{n}$\\\hline
Variance $\sigma^{2}$ & $s^{2}=\dfrac{1}{n-1}\left[\displaystyle\sum x^{2}-\dfrac{\left(\sum x\right)^{2}}{n}\right]$\newline or \ $s^{2}=\dfrac{1}{n-1}\displaystyle\sum(x-\xb)^{2}$ & $s^{2}=\dfrac{1}{n-1}\left[\displaystyle\sum(x-a)^{2}-\dfrac{\left(\sum(x-a)\right)^{2}}{n}\right]$\\\hline
\end{tabular}\end{center}
\gcs{\gckey{STAT} $\to$ \gc{CALC} $\to$ \gc{1-Var Stats}. Use \gc{Sx} ($=s$, divides by $n-1$), so $s^{2}=\gc{Sx}^{2}$. For a frequency table, enter the values in \gc{L1} and the frequencies in \gc{L2}, then \gc{1-Var Stats L1,L2}. Do \textbf{not} use \gc{$\sigma$x} (divides by $n$; this is not the unbiased estimate).}
\noindent\begin{minipage}[t]{0.48\textwidth}
\exlab $n=10$, $\sum x=125$, $\sum x^{2}=1830$
\cdisp{\xb&=\dfrac{125}{10}=12.5\\s^{2}&=\dfrac19\left[1830-\dfrac{125^{2}}{10}\right]\\&=\dfrac{267.5}{9}=29.7}
\end{minipage}\hfill
\begin{minipage}[t]{0.48\textwidth}
\exlab given $\sum(x-10)=25$, $\sum(x-10)^{2}=330$
\cdisp{\xb&=10+\dfrac{25}{10}=12.5\\s^{2}&=\dfrac19\left[330-\dfrac{25^{2}}{10}\right]\\&=\dfrac{267.5}{9}=29.7}
\end{minipage}
\begin{theorem}[before skip=3mm,after skip=3mm]{Unbiased estimate of population variance}
\[s^{2}=\dfrac{n}{n-1}\,\sigma_{x}^{2}\]
where $\sigma_{x}^{2}=\dfrac{1}{n}\displaystyle\sum(x-\xb)^{2}$ is the \textbf{sample variance} (GC: \gc{$\sigma$x}$^{2}$).
\end{theorem}
\noindent\exlab (Hypothesis testing) \hfill{\small[VJC 2025 Prelim P2 Q12]}\par
\noindent A hospital introduces a new medication regimen to help patients recover from a viral infection more quickly. Previously, the average recovery time was 7.2 days.\par
\noindent\ldots\par
\noindent It was discovered that an intern at the hospital had made an error in recording the recovery times. A random sample of another 30 patients was taken and their recovery times under the new regimen was taken and recorded by a senior researcher. The sample standard deviation was found to be 1.7 days and the sample mean time was denoted by $k$.\par
\noindent(f) A test at the 1\% significance level concludes that the average recovery time did not improve from 7.2 days. Find the set of values of $k$.\hfill[3]
\cdisp{s^{2}&=\dfrac{n}{n-1}(\text{sample variance})\\&=\dfrac{30}{29}(1.7)^{2}\\&=2.9897}
\clearpage
\noindent\exlab The masses, in kg, of a random sample of $5$ parcels are $4,\ 7,\ 8,\ 10,\ 11$. Find unbiased estimates of the population mean and variance.
\cdisp{\xb&=\dfrac{4+7+8+10+11}{5}=8\\\sum(x-\xb)^{2}&=(-4)^{2}+(-1)^{2}+0^{2}+2^{2}+3^{2}\\&=30\\s^{2}&=\dfrac{1}{n-1}\sum(x-\xb)^{2}\\&=\dfrac{30}{4}\\&=7.5}
\section{Distribution of the Sample Mean}
\begin{center}\bba\small
\begin{tabular}{|L{5.0cm}|M{4.6cm}|L{5.6cm}|}\hline
\textbf{Population} & \textbf{Distribution of $\Xb$} & \textbf{Reason}\\\Xhline{2pt}
Normal, $X\sim\No{\mu}{\sigma^{2}}$ & $\Xb\sim\No{\mu}{\dfrac{\sigma^{2}}{n}}$ \ (exact) & any $n$; linear combination of normals\\\hline
Not normal / unknown, $n$ \textbf{large} (e.g.\ $n\geqslant30$) & $\Xb\sim\No{\mu}{\dfrac{\sigma^{2}}{n}}$ approximately by CLT & \textbf{Central Limit Theorem}\\\hline
\end{tabular}\end{center}
\noindent\begin{minipage}[c]{0.46\textwidth}
\begin{tikzpicture}
\begin{axis}[nc,width=\linewidth,height=2.4cm,domain=-3.4:3.4]
\addplot[black,semithick]{\nf{0}};
\addplot[blue!60!black,semithick,samples=140]{3*exp(-0.5*((3*x)^2))/2.5066};
\draw[dashed,thin] (axis cs:0,0)--(axis cs:0,1.2);
\node[below] at (axis cs:0,0){$\mu$};
\node[font=\small,anchor=west] at (axis cs:1.2,0.34){$X$};\node[font=\small,blue!60!black,anchor=west] at (axis cs:0.45,0.9){$\Xb$};
\end{axis}\end{tikzpicture}
\end{minipage}\hfill
\begin{minipage}[c]{0.50\textwidth}\small
Same centre $\mu$. \ $\Xb$ is \textbf{narrower}: s.d.\ $\dfrac{\sigma}{\sqrt n}$ (larger $n$, narrower).
\exlab population mean $80$, s.d.\ $12$, $n=40$
\cdisp{\Xb&\sim\No{80}{\dfrac{12^{2}}{40}}\text{ approximately by CLT}\\\Pp{\Xb>83}&=0.0569}
\end{minipage}
\begin{theorem}[before skip=3mm,after skip=3mm]{Central Limit Theorem}
For a random sample of size $n$ (large, e.g.\ $n\geqslant30$) from \textbf{any} population with mean $\mu$ and variance $\sigma^{2}$:
\begin{itemize}[leftmargin=*,itemsep=1pt,topsep=1pt]
\item $\Xb\sim\No{\mu}{\dfrac{\sigma^{2}}{n}}$ approximately by CLT.
\item $X_{1}+X_{2}+\cdots+X_{n}\sim\No{n\mu}{n\sigma^{2}}$ approximately by CLT.
\end{itemize}
\end{theorem}
\clearpage
\section{Which Distribution?}
\begin{center}\ba\small
\begin{tabular}{|L{6.3cm}|L{3.4cm}|L{5.5cm}|}\hline
\textbf{The question \ldots} & \textbf{Use} & \textbf{Key step}\\\Xhline{2pt}
number of successes in $n$ independent trials & $\Bi{n}{p}$ & state the conditions\\\hline
one measurement, normal & $\No{\mu}{\sigma^{2}}$ & standardise / GC\\\hline
total or difference of independent normals & $\No{\cdot}{\cdot}$ & add the \textbf{variances}\\\hline
mean of $n$ values, population normal & $\Xb\sim\No{\mu}{\sigma^{2}/n}$ & exact\\\hline
mean of $n$ values, population not normal, $n$ large & $\Xb\sim\No{\mu}{\sigma^{2}/n}$ approximately by CLT & say ``by CLT''\\\hline
\end{tabular}\end{center}
\section{Further Question Types}
\begin{center}\renewcommand{\arraystretch}{1.12}\footnotesize
\begin{tabular}{|L{4.6cm}|L{10.6cm}|}\hline
\textbf{Type} & \textbf{Method}\\\Xhline{2pt}
Population given by a table or $\Bi{n}{p}$, $n$ observations large & $\mu=\E(X)$, $\sigma^{2}=\Var(X)$ first (Ch 3, 4); then $\Xb\sim\No{\mu}{\sigma^{2}/n}$ approximately by CLT\\\hline
Total of many independent counts (e.g.\ defects in $100$ batches) & $S=X_{1}+\cdots+X_{n}\sim\No{n\mu}{n\sigma^{2}}$ approximately by CLT\\\hline
Least $n$ for $\Pp{\Xb<c}\geqslant p$ & $\dfrac{c-\mu}{\sigma/\sqrt{n}}\geqslant$ \gc{invNorm(p)}, solve for $\sqrt n$; or GC table: $\mathrm{Y}_{1}=$ \gc{normalcdf} with s.d.\ $\sigma/\sqrt{\mathrm{X}}$, X as $n$\\\hline
Is the sample size big enough? & $n<30$ and population not known to be normal: CLT \textbf{cannot} be used; $n\geqslant30$: say ``by CLT''\\\hline
\end{tabular}\end{center}
\noindent\begin{minipage}[t]{0.48\textwidth}
\exlab The number of defective items, $X$, in a box of $12$ items is modelled by $X\sim\Bi{12}{0.25}$. A random sample of $40$ boxes is taken. Find the probability that the mean number of defective items per box exceeds $3.4$.
\cdisp{\E(X)&=12(0.25)=3\\\Var(X)&=12(0.25)(0.75)=2.25\\\because\ n&=40\text{ is large}\\\Xb&\sim\No{3}{\dfrac{2.25}{40}}\\&\quad\text{approximately by CLT}\\\Pp{\Xb>3.4}&=0.0458}
\end{minipage}\hfill
\begin{minipage}[t]{0.48\textwidth}
\exlab The number of faults, $X$, recorded by a sensor in an hour has the following probability distribution.
\begin{center}\ba\small\begin{tabular}{|c|c|c|c|c|}\hline
$x$&0&1&2&3\\\hline
$\Pp{X=x}$&0.4&0.3&0.2&0.1\\\hline
\end{tabular}\end{center}
$45$ independent hourly readings are taken. Find the probability that the total number of faults recorded exceeds $50$.
\cdisp{\E(X)&=0(0.4)+1(0.3)+2(0.2)+3(0.1)\\&=1\\\E(X^{2})&=0(0.4)+1(0.3)+4(0.2)+9(0.1)\\&=2\\\Var(X)&=2-1^{2}=1\\\text{let }S&=X_{1}+X_{2}+\cdots+X_{45}\\\because\ n&=45\text{ is large}\\S&\sim\No{45(1)}{45(1)}\\&\quad\text{approximately by CLT}\\S&\sim\No{45}{45}\\&\quad\text{approximately by CLT}\\\Pp{S>50}&=0.228}
\end{minipage}
